\section{A Comparative Map of Windows}

\begin{center}
\begin{tabularx}{\textwidth}{@{}l l >{\raggedright\arraybackslash}X >{\raggedright\arraybackslash}X@{}}
\toprule
Window type & Typical duration & Opening mechanism & Strategic use \\
\midrule
Morning reset & 5--20 minutes & Default inputs are not yet fully loaded and the prior day's frame has weakened & Best for cognitively heavy starts when the first task is already written down \\
Task boundary & 1--5 minutes & Switching cost is already being paid because one block has just ended & Best for chaining prepared work while the system is already in motion \\
Spatial transition & Variable & Physical relocation suppresses some old cues and makes a new role easier to name & Best for cue replacement and environment-led redirection \\
Social trigger & Variable & External commitment narrows admissible behaviour and raises the cost of drift & Best for actions that benefit from accountability or public start conditions \\
\bottomrule
\end{tabularx}
\end{center}

\begin{discussion}
Operationally, the rule is simple: identify the window, pre-position the first
action, remove the most likely closure trigger, and convert the cheap interval
into an irreversible start before inertia rebuilds. Morning resets are usually
best for deep work, task boundaries for chained execution, spatial transitions
for cue replacement, and social triggers for commitment-sensitive behaviours.
\end{discussion}
